\expandafter\ifx\csname CAPETINCLUDEMODE\endcsname\relax
\documentclass[11pt,a4paper]{article}
% preamble.tex — packages et commandes communes a toutes les fiches CAPET SII IE
\input{glyphtounicode}
\pdfgentounicode=1
\usepackage{cmap}
\usepackage[T1]{fontenc}
\usepackage[utf8]{inputenc}
\usepackage[french]{babel}
\usepackage[a4paper, top=1.6cm, bottom=1.6cm, left=1.4cm, right=1.4cm]{geometry}
\setlength{\headheight}{14.5pt}
\usepackage{amsmath, amssymb, amsfonts}
\usepackage{siunitx}
\sisetup{locale=FR, output-decimal-marker={,}, inter-unit-product=\ensuremath{{}\cdot{}}}
\usepackage{array, booktabs, multirow, tabularx}
\usepackage{xcolor}
\usepackage[most, breakable]{tcolorbox}
\usepackage{multicol}
\usepackage{enumitem}
\usepackage{mathtools}
\usepackage{lmodern}
\usepackage{microtype}
\usepackage{fancyhdr}
\usepackage{tikz}
\usetikzlibrary{positioning,arrows.meta}
\usepackage{lastpage}

% ── Couleurs par domaine ──────────────────────────────────────────────
\definecolor{colElec}{RGB}{26, 82, 148}       % Électrotechnique — bleu foncé
\definecolor{colElecLight}{RGB}{210, 225, 245} % fond bleu clair
\definecolor{colAuto}{RGB}{34, 120, 66}        % Automatique — vert
\definecolor{colMeca}{RGB}{185, 90, 20}        % Mécanique — orange
\definecolor{colTherm}{RGB}{160, 40, 40}       % Thermique — rouge
\definecolor{colQT}{RGB}{232, 244, 233}        % fond question type (vert très léger)
\definecolor{colJury}{RGB}{255, 243, 205}      % fond jury (jaune ambre)
\definecolor{colJuryFrame}{RGB}{200, 130, 0}   % bordure jury
\definecolor{colGrayLight}{RGB}{245, 245, 247} % fond gris neutre
\definecolor{colSeparator}{RGB}{180, 180, 180} % gris séparateur

% ── Style tcolorbox de section principale ────────────────────────────
\tcbset{
  ficheSection/.style={
    breakable,
    enhanced,
    colback=#1!7!white,
    colframe=#1,
    colbacktitle=#1,
    coltitle=white,
    fonttitle=\bfseries\normalsize,
    left=4pt, right=4pt, top=3pt, bottom=4pt,
    before skip=6pt, after skip=6pt,
  }
}

% ── Style encadré question type ──────────────────────────────────────
\tcbset{
  questiontype/.style={
    breakable,
    enhanced,
    colback=colQT,
    colframe=colElec!55!black,
    colbacktitle=colElec!55!black,
    coltitle=white,
    fonttitle=\bfseries\small,
    left=4pt, right=4pt, top=3pt, bottom=3pt,
    before skip=5pt, after skip=5pt,
  }
}

% ── Style encadré jury ───────────────────────────────────────────────
\tcbset{
  jury/.style={
    breakable,
    enhanced,
    colback=colJury,
    colframe=colJuryFrame,
    fonttitle=\bfseries\normalsize,
    left=4pt, right=4pt, top=3pt, bottom=3pt,
  }
}

% ── Mise en page ─────────────────────────────────────────────────────
\pagestyle{fancy}
\fancyhf{}
\renewcommand{\headrulewidth}{0.4pt}
\renewcommand{\footrulewidth}{0.3pt}

\setlist[itemize]{leftmargin=1.2em, itemsep=1pt, topsep=2pt, parsep=0pt}
\setlist[enumerate]{leftmargin=1.2em, itemsep=3pt, topsep=2pt, parsep=0pt}

\setlength{\parskip}{2pt}
\setlength{\columnsep}{10pt}

\def\CAPETBEGINFILE{\begin{document}}
\def\CAPETENDFILE{\end{document}}
\else
\def\CAPETBEGINFILE{}
\def\CAPETENDFILE{}
\fi

\newcommand{\totalPages}{5}  % mis - jour apres premiere compile

\fancyhead[L]{\textbf{\textcolor{colElec}{T02}} - Puissance en régime AC}
\fancyhead[R]{\small\textit{CAPET SII IE - Session 2026}}
\fancyfoot[C]{\small\thepage\,/\,\pageref{LastPage}}

\CAPETBEGINFILE

% -----------------------------------------------------------------------
%  BLOC EN-TETE (integre - la premiere page, sans saut de page)
% -----------------------------------------------------------------------

\begin{tcolorbox}[
  enhanced, breakable,
  colback=colElec,
  colframe=colElec,
  left=6pt, right=6pt, top=5pt, bottom=5pt,
  before skip=0pt, after skip=4pt
]
  {\color{white}\bfseries\large T02 \quad Puissance en régime alternatif sinusoïdal (AC)}\\[2pt]
  {\color{colElecLight}\small
    \textbf{Domaine :} électrotechnique \quad
    \textbf{Présence :} \textbf{13/13} sujets (2013-2026) - theme \emph{systématiquement} présent
  }\\[1pt]
  {\color{colElecLight!80!white}\footnotesize
    Sous-thèmes : triangle des puissances - facteur de puissance - triphasé équilibré - compensation réactive - bilan de puissance
  }
\end{tcolorbox}

% -----------------------------------------------------------------------
%  1. RAPPELS DE COURS
% -----------------------------------------------------------------------

\begin{tcolorbox}[
  ficheSection=colElec,
  title={1. Rappels de cours},
]

% -- 1.1 Monophase --------------------------------------------------
\textbf{1.1 Régime sinusoïdal monophasé}

\smallskip
Récepteur d'impédance $\underline{Z} = R + jX$, alimenté par $u(t) = U\sqrt{2}\cos(\omega t)$.
Courant : $i(t) = I\sqrt{2}\cos(\omega t - \varphi)$, avec $\varphi = \arg(\underline{Z})$.

\smallskip
\begin{center}
\renewcommand{\arraystretch}{1.45}
\begin{tabular}{llll}
\toprule
\textbf{Grandeur} & \textbf{Symb.} & \textbf{Expression} & \textbf{Unité} \\
\midrule
Puissance active    & $P$   & $P = UI\cos\varphi = RI^2 = \dfrac{U^2\cos\varphi}{Z}$ & W \\[5pt]
Puissance réactive  & $Q$   & $Q = UI\sin\varphi = XI^2 = \dfrac{U^2\sin\varphi}{Z}$ & VAR \\[5pt]
Puissance apparente & $S$   & $S = UI = \sqrt{P^2+Q^2}$ & VA \\[4pt]
Facteur de puissance & $\cos\varphi$ & $\cos\varphi = \dfrac{P}{S} = \dfrac{R}{Z}$ & - $\in[0;1]$ \\[5pt]
Puissance complexe  & $\underline{S}$ & $\underline{S} = \underline{U}\,\overline{\underline{I}} = P + jQ$ & VA \\
\bottomrule
\end{tabular}
\end{center}

\smallskip
\textbf{Triangle des puissances :}

\begin{center}
\begin{minipage}[c]{0.32\textwidth}
% Schema TikZ
\begin{tikzpicture}[scale=1.05, font=\small]
  \draw[thick, colElec, ->] (0,0) -- (3.2,0) node[right] {$P$};
  \draw[thick, colElec!60!black, ->] (3.2,0) -- (3.2,1.8) node[right] {$Q$};
  \draw[thick, colElec!40!black, ->] (0,0) -- (3.2,1.8) node[above left=-2pt] {$S$};
  % angle phi
  \draw[colElec!70!black] (0.7,0) arc[start angle=0, end angle=29.3, radius=0.7];
  \node at (1.05, 0.2) {\footnotesize$\varphi$};
  % labels
  \node[below] at (1.6,0) {$P = S\cos\varphi$};
  \node[right] at (3.25,0.9) {$Q = S\sin\varphi$};
\end{tikzpicture}
\end{minipage}
\hfill
\begin{minipage}[c]{0.30\textwidth}
\[S^2 = P^2 + Q^2\]
\[\tan\varphi = \frac{Q}{P}\]
\end{minipage}
\hfill
\begin{minipage}[c]{0.30\textwidth}
\footnotesize
$Q > 0$ : inductif (moteur)\\
$Q < 0$ : capacitif (condo)\\
$Q = 0$ : resistif, $\cos\varphi=1$
\end{minipage}
\end{center}

\medskip
\hrule\vspace{4pt}

% -- 1.2 Triphase --------------------------------------------------
\textbf{1.2 Régime triphasé équilibré}

$V$ = tension simple (phase-neutre), $U$ = tension composée (entre phases) :
$\displaystyle U = \sqrt{3}\,V$

\smallskip
\begin{center}
\renewcommand{\arraystretch}{1.4}
\begin{tabular}{lll}
\toprule
\textbf{Grandeur} & \textbf{Expression} & \textbf{Remarque} \\
\midrule
Puissance active    & $P = 3VI\cos\varphi = \sqrt{3}\,UI\cos\varphi$ & $I$ = courant de ligne \\[4pt]
Puissance réactive  & $Q = 3VI\sin\varphi = \sqrt{3}\,UI\sin\varphi$ & \\[4pt]
Puissance apparente & $S = 3VI = \sqrt{3}\,UI = \sqrt{P^2+Q^2}$ & \\[4pt]
Facteur de puissance & $\cos\varphi = P/S$ & identique au monophasé \\
\bottomrule
\end{tabular}
\end{center}

\smallskip
\textbf{Couplage étoile (Y) / triangle ($\Delta$) :}

\smallskip
\begin{center}
\renewcommand{\arraystretch}{1.25}
\begin{tabular}{lcc}
\toprule
 & \textbf{étoile (Y)} & \textbf{Triangle ($\Delta$)} \\
\midrule
Tension aux bornes du recepteur & $V = U/\sqrt{3}$ & $U$ \\
Courant dans le recepteur & $I_\text{ligne}$ & $I_\text{ligne}/\sqrt{3}$ \\
\bottomrule
\end{tabular}
\end{center}

\medskip
\hrule\vspace{4pt}

% -- 1.3 Compensation ----------------------------------------------
\textbf{1.3 Compensation de l'énergie réactive}

Objectif : rôlever $\cos\varphi_1$ - $\cos\varphi_2 > \cos\varphi_1$ en ajoutant des condensateurs en parallèle.

\[
  Q_C = P\,\bigl(\tan\varphi_1 - \tan\varphi_2\bigr)
  \qquad \text{[puissance réactive - fournir par les condensateurs]}
\]

\begin{center}
\renewcommand{\arraystretch}{1.3}
\begin{tabular}{lll}
\toprule
\textbf{Couplage} & \textbf{Tension sur le condo} & \textbf{Capacite par phase} \\
\midrule
étoile (Y) & $V = U/\sqrt{3}$ & $\displaystyle C = \dfrac{Q_C}{3\,\omega\,V^2}$ \\[6pt]
Triangle ($\Delta$) & $U$ & $\displaystyle C = \dfrac{Q_C}{3\,\omega\,U^2}$ \\[4pt]
\bottomrule
\end{tabular}
\end{center}
\vspace{2pt}
{\footnotesize Sur réseau 400\,V/230\,V : $V=230$\,V, $\omega = 2\pi\times50 \approx 314$\,rad/s. Un condensateur $\Delta$ vaut 3 fois moins qu'un equivalent Y.}

\medskip
\hrule\vspace{4pt}

% -- 1.4 Bilan installation ----------------------------------------
\textbf{1.4 Bilan de puissance d'une installation complexe}

Plusieurs recepteurs en parallèle - \textbf{sommer $P$ et $Q$ séparément, jamais $S$} :
\[
  P_\text{tot} = \textstyle\sum P_k \qquad Q_\text{tot} = \textstyle\sum Q_k \qquad S_\text{tot} = \sqrt{P_\text{tot}^2 + Q_\text{tot}^2}
\]
Pour un moteur asynchrone ($P_u$ = puissance utile mecanique, $\eta$ = rendement) :
\[
  P_a = \frac{P_u}{\eta} \qquad
  Q = P_a\tan\varphi \qquad
  S = \frac{P_a}{\cos\varphi}
\]

\end{tcolorbox}

% -----------------------------------------------------------------------
%  2. FORMULAIRE SYNTHeTIQUE (2 colonnes)
% -----------------------------------------------------------------------

\begin{tcolorbox}[
  ficheSection=colElec,
  title={2. Formulaire synthétique},
]

\small
\begin{multicols}{2}
\renewcommand{\arraystretch}{1.45}
\begin{tabular}{@{}p{3.8cm}ll@{}}
\toprule
\textbf{Formule} & \textbf{Usage} & \textbf{U.} \\
\midrule
\multicolumn{3}{@{}l}{\footnotesize\textit{Monophase}} \\
$P = UI\cos\varphi$ & P active & W \\
$Q = UI\sin\varphi$ & P réactive & VAR \\
$S = \sqrt{P^2+Q^2}$ & P apparente & VA \\
$\cos\varphi = P/S = R/Z$ & Facteur puis. & - \\
\midrule
\multicolumn{3}{@{}l}{\footnotesize\textit{Triphase équilibré}} \\
$P = \sqrt{3}\,UI\cos\varphi$ & P active & W \\
$Q = \sqrt{3}\,UI\sin\varphi$ & P réactive & VAR \\
$S = \sqrt{3}\,UI$ & P apparente & VA \\
$U = \sqrt{3}\,V$ & Tension & V \\
$I = P/(\sqrt{3}\,U\cos\varphi)$ & Courant ligne & A \\
\bottomrule
\end{tabular}

\columnbreak

\renewcommand{\arraystretch}{1.45}
\begin{tabular}{@{}p{3.8cm}ll@{}}
\toprule
\textbf{Formule} & \textbf{Usage} & \textbf{U.} \\
\midrule
\multicolumn{3}{@{}l}{\footnotesize\textit{Compensation réactive}} \\
$Q_C = P(\tan\varphi_1 - \tan\varphi_2)$ & - compenser & VAR \\[2pt]
$C = Q_C/(3\omega V^2)$ & Condo étoile & F \\[2pt]
$C = Q_C/(3\omega U^2)$ & Condo triangle & F \\
\midrule
\multicolumn{3}{@{}l}{\footnotesize\textit{Moteur asynchrone}} \\
$P_a = P_u / \eta$ & P absorbée & W \\
$Q = P_a \tan\varphi$ & P réactive & VAR \\
$S = P_a / \cos\varphi$ & P apparente & VA \\
\midrule
\multicolumn{3}{@{}l}{\footnotesize\textit{Trigonometrie}} \\
$\sin\varphi = \sqrt{1-\cos^2\!\varphi}$ & & - \\
$\tan\varphi = \sin\varphi/\cos\varphi$ & & - \\
\bottomrule
\end{tabular}

\end{multicols}
\normalsize

\end{tcolorbox}

% -----------------------------------------------------------------------
%  3. QUESTIONS TYPES
% -----------------------------------------------------------------------

\begin{tcolorbox}[
  ficheSection=colElec,
  title={3. Questions types - extraites des sujets réels},
]

% -- QT1 -------------------------------------------------------------

% Mini-procedure de resolution type copie concours
\begin{tcolorbox}[questiontype,title={Procedure type copie concours (5 étapes)}]
\begin{enumerate}[leftmargin=1.5em]
  \item \textbf{Qualifier le modele} : hypotheses, domaine de validite, grandeurs d'entree/sortie.
  \item \textbf{Ecrire les lois} : equations de base avec symboles et unités SI.
  \item \textbf{Isoler l'inconnue} : expression litterale, puis application numerique tracee.
  \item \textbf{Contrôler l'ordre de grandeur} : coherence physique et unite finale.
  \item \textbf{Conclure concours} : phrase courte orientee decision technique.
\end{enumerate}
\end{tcolorbox}

% Visuel de methode : chemin de resolution CAPET
\begin{center}
\begin{tikzpicture}[font=\small, node distance=7mm and 6mm, >=Stealth]
  \node[draw, rounded corners, fill=colElec!8, inner sep=4pt] (n1) {Notions};
  \node[draw, rounded corners, fill=colElec!8, right=of n1, inner sep=4pt] (n2) {Méthode};
  \node[draw, rounded corners, fill=colElec!8, right=of n2, inner sep=4pt] (n3) {Application annale};
  \node[draw, rounded corners, fill=colElec!8, right=of n3, inner sep=4pt] (n4) {Vérification jury};
  \draw[->, thick, colElec] (n1) -- (n2);
  \draw[->, thick, colElec] (n2) -- (n3);
  \draw[->, thick, colElec] (n3) -- (n4);
\end{tikzpicture}
\end{center}
\begin{tcolorbox}[questiontype,
  title={QT1 - Bilan de puissance d'une installation triphasée \hfill\normalfont\footnotesize E2 disciplinaire 2024}
]

\textbf{énoncé.}
Un atelier triphasé $U=\SI{400}{\volt}$, $f=\SI{50}{\hertz}$ comporte :
un moto-extracteur ($P_u=\SI{100}{\kilo\watt}$, $\cos\varphi_\text{mot}=0{,}85$, $\eta=0{,}92$) ;
des ordinateurs ($P_0=\SI{1}{\kilo\watt}$, $\cos\varphi_0=1$) ;
eclairage et équipements ($P_e=\SI{43}{\kilo\watt}$, $Q_e=\SI{25{,}5}{\kilo VAR}$).
\textbf{Calculer} $P_\text{ch}$, $Q_\text{ch}$, $S_\text{ch}$, $I_\text{abs}$ et $F_p$.

\tcbline

\textbf{Corrigé.}

\textbf{étape 1 - puissances du moteur.}
\[
  P_\text{mot} = \frac{P_u}{\eta} = \frac{100}{0{,}92} \approx \SI{108{,}7}{\kilo\watt}
  \qquad
  \sin\varphi_\text{mot} = \sqrt{1-0{,}85^2} = 0{,}527
\]
\[
  Q_\text{mot} = P_\text{mot}\,\tan\varphi_\text{mot} = 108{,}7 \times \frac{0{,}527}{0{,}85} \approx \SI{67{,}4}{\kilo VAR}
\]

\textbf{étape 2 - bilan global} (sommer $P$ et $Q$ séparément, jamais $S$).
\[
  P_\text{ch} = 108{,}7 + 1 + 43 = \SI{152{,}7}{\kilo\watt}
  \qquad
  Q_\text{ch} = 67{,}4 + 0 + 25{,}5 = \SI{92{,}9}{\kilo VAR}
\]
\[
  S_\text{ch} = \sqrt{P_\text{ch}^2+Q_\text{ch}^2} = \sqrt{152{,}7^2+92{,}9^2} \approx \SI{179{,}1}{\kilo VA}
\]

\textbf{étape 3 - courant et facteur de puissance.}
\[
  I_\text{abs} = \frac{S_\text{ch}}{\sqrt{3}\,U} = \frac{179\,100}{\sqrt{3}\times 400} \approx \SI{258{,}5}{\ampere}
  \qquad
  F_p = \frac{P_\text{ch}}{S_\text{ch}} = \frac{152{,}7}{179{,}1} \approx \mathbf{0{,}853}
\]

\end{tcolorbox}

\vspace{4pt}

% -- QT2 -------------------------------------------------------------

\begin{tcolorbox}[questiontype,
  title={QT2 - Compensation de l'énergie réactive par condensateurs \hfill\normalfont\footnotesize E2 disciplinaire 2024}
]

\textbf{énoncé.}
(Donnees jury : $P=\SI{178{,}3}{\kilo\watt}$, $Q=\SI{108{,}7}{\kilo VAR}$.)
Un gestionnaire impose $F_p=1$ via une batterie de condensateurs étoile ($C$ par phase).
\textbf{B3.} Montrer que $C = Q_\text{abs}/(3\omega V_n^2)$.
\textbf{B4.} Calculer $C$.

\tcbline

\textbf{Corrigé.}

\textbf{B3.}
Pour $F_p'=1$ : $Q_\text{total} = Q_\text{abs} + Q_C = 0$, donc $Q_C = -Q_\text{abs}$.
Un condensateur étoile sous tension simple $V_n$ fournit $Q_C = 3C\omega V_n^2$, d'où :
\[
  \boxed{C = \frac{Q_\text{abs}}{3\,\omega\,V_n^2}}
\]

\textbf{B4.}
$V_n = \SI{230}{\volt}$, $\omega = 314$\,rad/s :
\[
  C = \frac{108\,700}{3 \times 314 \times 230^2} = \frac{108\,700}{49\,858\,200} \approx \SI{2{,}18}{\milli\farad}
\]

\end{tcolorbox}

\vspace{4pt}

% -- QT3 -------------------------------------------------------------

\begin{tcolorbox}[questiontype,
  title={QT3 - Relevement du facteur de puissance vers une valeur cible \hfill\normalfont\footnotesize E2 2017-2020}
]

\textbf{énoncé.}
Installation triphasée : $P=\SI{120}{\kilo\watt}$, $\cos\varphi_1=0{,}75$.
On souhaite $\cos\varphi_2=0{,}93$.
Calculer $Q_C$ et $C$ (couplage étoile, $V=\SI{230}{\volt}$).

\tcbline

\textbf{Corrigé.}
\[
  \tan\varphi_1 = \frac{\sqrt{1-0{,}75^2}}{0{,}75} = 0{,}882
  \qquad
  \tan\varphi_2 = \frac{\sqrt{1-0{,}93^2}}{0{,}93} = 0{,}396
\]
\[
  Q_C = P(\tan\varphi_1-\tan\varphi_2) = 120\times(0{,}882-0{,}396) = \SI{58{,}3}{\kilo VAR}
\]
\[
  C = \frac{58\,300}{3\times 314\times 230^2} \approx \SI{1{,}17}{\milli\farad}
\]
\textit{Vérif. :}
\[
  \cos\varphi_2 = \frac{120}{\sqrt{120^2+(105{,}8-58{,}3)^2}} = \frac{120}{129{,}0} \approx 0{,}930 \;\checkmark
\]

\end{tcolorbox}

\end{tcolorbox}

% -----------------------------------------------------------------------
%  4. ATTENTES DU JURY
% -----------------------------------------------------------------------

\begin{tcolorbox}[
  jury,
  title={4. Attentes du jury - Erreurs récurrentes et points de vigilance},
]

\begin{enumerate}[label=\textbf{\arabic*.}, leftmargin=*, itemsep=5pt]

  \item \textbf{Ne jamais sommer les puissances apparentes $S_k$.}
  $S_\text{tot} \neq S_1+S_2+\cdots$ \textit{(erreur eliminatoire)}.
  Méthode : sommer $P_k$, sommer $Q_k$, puis $S=\sqrt{P^2+Q^2}$.

  \item \textbf{Confondre tension simple et composée.}
  $V = U/\sqrt{3}$ - réseau 400\,V/230\,V : $U=400$\,V entre phases, $V=230$\,V phase-neutre.
  Les formules $P=\sqrt{3}\,UI\cos\varphi$ utilisent $U$ composée et $I$ de ligne.

  \item \textbf{Oublier le rendement du moteur.}
  $P_\text{absorbée} = P_\text{utile}/\eta$. C'est $P_a$ qui sert - calculer $Q$ et $S$, pas $P_u$.

  \item \textbf{Erreurs d'unités et d'ordre de grandeur.}
  Le jury signale chaque année des erreurs de facteur $10^3$.
  Convertir systématiquement en SI avant toute application numerique.

  \item \textbf{Confondre couplage étoile et triangle pour $C$.}
  étoile : $C = Q_C/(3\omega V^2)$ avec $V$ simple.
  Triangle : $C = Q_C/(3\omega U^2)$ avec $U$ composée.
  Un condensateur $\Delta$ vaut 3 fois moins qu'un condensateur Y equivalent.

  \item \textbf{Signe de la puissance réactive.}
  Inductances ($Q>0$), condensateurs ($Q<0$).
  La compensation annule la somme algebrique des $Q_k$.

  \item \textbf{Rédiger la demonstration, pas seulement le resultat.}
  Exemple attendu : \textit{"Pour $F_p=1$, $Q_\text{tot}=Q_\text{abs}+Q_C=0$, donc $Q_C=-Q_\text{abs}$.
  Un condensateur étoile donne $Q_C=3C\omega V_n^2$, d'où $C=Q_\text{abs}/(3\omega V_n^2)$."}

  \item \textbf{Verifier le resultat.}
  Contrôler que $P\leq S$, $\cos\varphi\in[0;1]$, valeur de $C$ realiste (mF - eF selon les puissances).

\end{enumerate}

\medskip
\begin{tcolorbox}[
  colback=colJuryFrame!12!colJury,
  colframe=colJuryFrame,
  left=4pt, right=4pt, top=2pt, bottom=2pt,
  fontupper=\small\itshape
]
Rapports jury 2019-2025 : "La manipulation des unités et des puissances de 10 est tres souvent hasardeuse, entraînant des resultats dont l'ordre de grandeur meme est faux."
"Un nombre significatif de candidats ne parvient pas - caracteriser des constituants electriques ou - identifier leur rôle dans une chaîne de puissance."
\end{tcolorbox}

\end{tcolorbox}

% -----------------------------------------------------------------------
%  PIED - PRESENCE DANS LES SUJETS
% -----------------------------------------------------------------------

\begin{tcolorbox}[
  enhanced, breakable,
  colback=colGrayLight,
  colframe=colSeparator,
  left=4pt, right=4pt, top=3pt, bottom=3pt,
  before skip=4pt, after skip=0pt,
  fonttitle=\small\bfseries,
  title={Présence dans les sujets analyses (2013-2026)},
]
\begin{center}
\footnotesize\renewcommand{\arraystretch}{1.15}
\begin{tabular}{lccccccccccccc}
\toprule
\textbf{Annee} & 13 & 14 & 15 & 16 & 17 & 18 & 19 & 20 & 21 & 22 & 23 & 24 & 25 \\
\midrule
E1 (commune)    & $\checkmark$ & $\checkmark$ & $\checkmark$ & $\checkmark$ & $\checkmark$ & $\checkmark$ & $\checkmark$ & $\checkmark$ & $\checkmark$ & $\checkmark$ & $\checkmark$ & $\checkmark$ & $\checkmark$ \\
E2 discipl. IE  & - & - & - & - & $\checkmark$ & $\checkmark$ & $\checkmark$ & $\checkmark$ & $\checkmark$ & $\checkmark$ & $\checkmark$ & $\checkmark$ & $\checkmark$ \\
\bottomrule
\end{tabular}
\end{center}
\smallskip
\footnotesize\textbf{Sous-thèmes frequents :}
$P/Q/S$ triphasé (13/13) - facteur de puissance (13/13) - bilan installation (11/13) - compensation condensateurs (9/13) - chute de tension ligne (6/13)

\medskip
\textbf{Ordres de grandeur (repères rapides).}
\begin{itemize}[leftmargin=1.3em]
  \item Reseau BT FR : tension phase-neutre typique \SI{230}{\volt}, tension phase-phase typique \SI{400}{\volt}, fréquence \SI{50}{\hertz}.
  \item Courants d'etude frequents en sujet : de quelques \si{\ampere} a quelques dizaines de \si{\ampere}.
  \item En conversion d'énergie, un rendement attendu reste generalement entre 0{,}8 et 0{,}98 selon la technologie.
\end{itemize}

\medskip
\textbf{Synthese express (à refaire sans support).}
\begin{enumerate}[leftmargin=1.4em]
  \item Poser les hypotheses du modele et les conventions de signe.
  \item Ecrire 2 a 4 relations clés avec unités.
  \item Réaliser un calcul numerique complet et verifier l'ordre de grandeur.
  \item Justifier le resultat par un argument technique de niveau CAPET.
\end{enumerate}

\end{tcolorbox}

\CAPETENDFILE










